%!TEX root = ../hutbthesis_main.tex
% 设置中文摘要
\keywordscn{肌肉骨骼控制\quad 运动模仿\quad 强化学习\quad 量化评估\quad Kinesis\quad 生物合理性算法优化}
\begin{abstractzh}

现在的肌肉骨骼运动模仿算法有不少短板，比如模仿精度不高，动作不够稳定，还有生物合理性差。针对这些问题，我在Kinesis仿真平台上深入研究了更符合生物规律的肌肉骨骼运动控制算法。核心方法是用深度强化学习，同时设计了多目标奖励函数，把姿态误差、速度误差、生物约束和任务完成度这四个惩罚项引入进来，目的是让运动模仿更准确、动作更自然。评估标准选了三个：平均每关节位置误差（MPJPE）、帧覆盖率和任务成功率，分别衡量模仿精度、动作稳定性和任务鲁棒性。对比实验结果很直接：引入多目标奖励和生物约束惩罚项后，关节误差明显降低，动作崩溃情况减少，整体运动看起来更自然。换句话说，这套方法在肌肉骨骼运动控制上确实有效，而且表现优越。无论是康复机器人、虚拟人动画还是外骨骼控制，这项研究都能提供有价值的算法参考。

\end{abstractzh}
